\section{Methodology}

This study was conducted as a systematic literature review following Kitchenham (2007) \cite{kitchenham2007}. As a secondary study, it adopts a descriptive-analytical synthesis that combines quantitative summarization of the extracted evidence with narrative interpretation across studies on fuzzy-logic-based inference systems in real-time air-quality monitoring platforms.

\subsection{Study Design}
The study design is structured through the \textbf{PICo} model shown in Table~\ref{tab:pico_air_quality}, in order to delimit precisely the phenomenon under study, the methodological focus, and the application context.

\begin{longtblr}
[
 label={tab:pico_air_quality},
 caption={PICo framework used for the state-of-the-art review. Own elaboration.}]{
colspec={X[j,m,0.7] X[j,m,3] X[j,m,5]},
width=\textwidth,
rowhead=1,
row{1}={font=\bfseries}
}
\toprule
\textbf{Compo- nent} & \textbf{Condensed term} & \textbf{Operational description} \\
\midrule

P &
Urban air quality &
Evaluated from atmospheric pollutant concentrations (e.g., PM$_{2.5}$, PM$_{10}$, NO$_2$, O$_3$) and associated environmental variables used as indicators of air-quality status. \\

I &
Fuzzy logic for environmental assessment &
Application of fuzzy inference systems (e.g., Mamdani or Sugeno) to evaluate air-quality level and associated risk through linguistic rules and membership functions. \\

Co &
Real-time monitoring &
Real-time monitoring platforms and systems oriented toward the analysis and visualization of air-quality data from environmental sensors and/or open data sources. \\

\bottomrule
\end{longtblr}


This approach makes it possible to identify and systematically analyze existing contributions that combine air-quality assessment, fuzzy logic, and real-time monitoring platforms.

From the delimitation established through the PICo model, the following research question is formulated:

\textbf{Research question:}

\begin{quote}
How are fuzzy-logic-based inference systems used in real-time monitoring platforms to evaluate urban air quality and the associated risks?
\end{quote}

In order to operationalize this question and structure the review of the state of the art systematically, the following research directives are defined:

\begin{itemize}
    \item \textbf{Directive 1 (Data and variables):}
    Which atmospheric pollutants and environmental variables are most frequently used to evaluate urban air quality in real-time monitoring platforms?

    \item \textbf{Directive 2 (Fuzzy-logic approach):}
    Which types of fuzzy-logic-based inference systems are used to evaluate air quality and the associated risk, and what are their main design criteria?

    \item \textbf{Directive 3 (Platform and real time):}
    How are real-time monitoring platforms implemented for the analysis and visualization of air quality, and which architectures and technologies are used?
\end{itemize}

\subsection*{Search Strategy}
The literature search was carried out in Scopus, Web of Science, and IEEE Xplore. This combination was used to retrieve studies indexed across the engineering, computing, and environmental domains in which research on fuzzy-logic-based air-quality monitoring is typically published. Based on these sources, the search strategy was defined using the PICo framework established above, with the aim of identifying relevant literature on the use of fuzzy-logic-based inference systems in real-time air-quality monitoring platforms. The search was restricted to English-language publications from the 2021--2026 period and included peer-reviewed journal articles and conference proceedings.

In building the search string, the core concepts of each PICo component (Problem/Phenomenon, Interest, and Context) were identified, and variant terms and equivalent expressions recurrently used in the scientific literature were derived. Table~\ref{tab:derivacion_pico} presents the derived terms for each PICo component.

\begin{longtblr}
[
 label={tab:derivacion_pico},
 caption={Derivación de términos de búsqueda a partir del marco PICo. Elaboración propia.}]{
  colspec={l l X[l,m]},
  width=\textwidth,
  rowhead=1,
  row{1}={font=\bfseries}
}
\toprule
\textbf{Componente} & \textbf{Concepto nuclear} & \textbf{Términos derivados y variantes} \\
\midrule

P &
Calidad del aire &
air; ambient; atmospheric; quality; pollution; pollutant*; ``particulate matter''; PM$_{2.5}$; PM$_{10}$; NO$_2$; O$_3$; CO; SO$_2$ \\

I &
Lógica difusa &
fuzzy; ANFIS; ``neuro-fuzzy''; logic; inference; ``expert system''; rule; Mamdani; Sugeno; ``type-2 fuzzy''; ``interval type-2'' \\

Co &
Monitoreo en tiempo real &
``real-time''; ``real time''; monitoring; platform; system; sensor*; IoT; ``Internet of Things''; WSN \\

\bottomrule
\end{longtblr}


Based on these terms, one search block was constructed for each PICo component. Within each block, OR was used to group synonymous or related terms, whereas AND was used across the P, I, and Co blocks so that the retrieved records addressed the three components simultaneously. The search strings were then adjusted in each database through title and abstract screening to improve thematic fit and control the volume of retrieved records.

The final search string used in the selected databases is presented in Table~\ref{tab:cadena_final}.

\begin{longtblr}
[
 label={tab:cadena_final},
 caption={Final search string used in the selected databases. Own elaboration.}]{
  colspec={X[l,m]},
  width=\textwidth,
  rowhead=1,
  row{1}={font=\bfseries}
}
\toprule
Final search string \\
\midrule

\small
(air OR ambient OR atmospheric) AND (quality OR pollution OR pollutant* OR ``particulate matter'' OR PM2.5 OR PM10 OR NO2 OR O3 OR CO OR SO2) \\
AND \\
(fuzzy OR ANFIS OR ``neuro-fuzzy'') AND (logic OR inference OR ``expert system'' OR rule OR Mamdani OR Sugeno OR ``type-2 fuzzy'' OR ``interval type-2'') \\
AND \\
(``real-time'' OR ``real time'' OR monitoring) AND (monitoring OR platform OR system) AND (sensor* OR IoT OR ``Internet of Things'' OR WSN) \\

\bottomrule
\end{longtblr}


This formulation retrieved a set of studies that covered the topic while remaining focused enough for systematic analysis.

\subsection*{Article Selection}
Study selection followed the PRISMA 2020 sequence of identification, screening, full-text retrieval, eligibility assessment, and final inclusion \cite{prisma2020}. After consolidating the records retrieved from the selected databases and removing duplicates, the set used for title and abstract screening was defined. The operational inclusion and exclusion criteria applied during screening and eligibility are presented in Table~\ref{tab:criterios_ie}.

\begin{adjustwidth}{-\extralength}{0cm}
\begin{longtblr}[
  label={tab:criterios_ie},
  caption={Inclusion and exclusion criteria applied in the \textit{screening} and eligibility phases, aligned with PICo, research directives (D1--D3), and Kitchenham's practical criteria. Own elaboration.}]
{
    colspec = {X[m,j,3] X[m,j,5] X[m,j,5] X[m,j,1.25]},
    width=\fulllength,
    rowhead=1,
    row{1}={font=\bfseries}
}
\toprule
\textbf{Criterion} &
\textbf{Inclusion (Operational description)} &
\textbf{Exclusion (Operational description)} &
\textbf{Phase} \\
\midrule

Relation to the phenomenon (Air quality -- P) &
The study explicitly addresses air quality or ambient atmospheric pollution. &
The study focuses on water, soil, agriculture, aquaculture, noise, or another domain unrelated to air quality. &
Screening \\

Central use of fuzzy logic (I) &
The title or abstract mentions \textit{fuzzy logic} or fuzzy inference systems as the main approach. &
No use of fuzzy logic is identified, or it does not constitute the study's central approach. &
Screening \\

Air-quality evaluation through fuzzy logic (P + I) &
The study proposes evaluation, classification, or analysis of air-quality level using fuzzy logic. &
Fuzzy logic is applied only to control (PID, HVAC, drying, ventilation, furnaces, etc.) without evaluating air quality as the main variable. &
Screening \\

Nature of the study (Article type) &
Primary study with an original technical proposal or implementation. &
Systematic review, narrative review, \textit{mapping study}, \textit{proceedings} volume, or grey literature. &
Screening \\

Air-quality monitoring platform or system (Co) &
The study proposes or implements a platform, monitoring system, IoT architecture, sensor integration, or structured use of \textit{datasets} for air-quality analysis. &
The study does not describe a system, platform, or data integration, or is limited to theoretical analysis without a monitoring or structured processing component. &
Screening \\

\textbf{(D1) }Specification of variables and pollutants used &
The study explicitly identifies the atmospheric pollutants and/or environmental variables used as inputs (e.g., PM$_{2.5}$, PM$_{10}$, NO$_2$, O$_3$, CO, SO$_2$, temperature, humidity, wind). &
The variables used are not clearly specified, or they cannot be identified from the full text. &
Full Text \\

\textbf{(D2) }Explicit application of a fuzzy inference system to air quality &
Fuzzy inference constitutes the central mechanism for evaluating air quality and/or associated risk. &
Fuzzy logic is tangential, secondary, or not directly applied to air-quality evaluation. &
Full Text \\

\textbf{(D2) }Structural description of the fuzzy model (type and design) &
The system type (e.g., Mamdani, Sugeno, ANFIS, type-2, or other) is specified, and structural elements (rules, membership functions, or architecture) are described. &
Neither the fuzzy-system type nor its structure is detailed, preventing characterization and systematic extraction. &
Full Text \\

\textbf{(D3) }Implementation in a monitoring system or platform &
The study describes implementation in a platform, IoT architecture, monitoring system, or structured air-quality data-processing environment. &
The model is presented without the context of a monitoring system/platform or without operational integration into a data-processing environment. &
Full Text \\

\textbf{(D3) }Description of application setting and unit of analysis (Setting / Participants) &
The application context (e.g., urban/city, stations, sensors, \textit{datasets}) and the evaluated unit of analysis are specified. &
The application setting or unit of analysis is not clearly described, preventing comparison with other studies. &
Full Text \\

\textbf{(Cross-cutting) }Documented methodological design and validation (Research Design / Sampling) &
The study describes an experimental or applied design and presents empirical validation or quantitative evaluation (e.g., metrics, comparison, test cases). &
It does not present a clear methodological design and/or does not report empirical validation or quantitative evaluation of the proposed system. &
Full Text \\

\bottomrule
\end{longtblr}
\end{adjustwidth}


During title and abstract screening, each record was assessed in terms of its relationship to the air-quality phenomenon (P), the central use of fuzzy logic (I), air-quality evaluation through fuzzy logic (P+I), the nature of the study, and the existence of a monitoring platform or system (Co). This phase made it possible to separate excluded studies from those that had to proceed to full-text retrieval.

Full-text retrieval combined an automated strategy and a manual strategy. In the automated phase, each record with a DOI was first queried in Unpaywall and, when no direct PDF was identified, a second query was made in OpenAlex to locate open-access links. This stage enabled traceable classification of downloaded cases and non-retrievable cases due to missing DOI, lack of a direct PDF link, or access restriction. For cases not obtained automatically, a manual search was conducted in the indexed database platforms, arXiv, and ResearchGate; when the document was not directly accessible, a request was managed through ResearchGate. Retrieval status was documented in the screening matrix through the \texttt{Retrieval\_result} field.

Eligibility was assessed on the retrieved full-text studies by applying criteria D1--D3 and the cross-cutting criterion of methodological design and validation, with explicit page-based evidence recording. At this stage, a single exclusion reason was assigned to each study when applicable, maintaining consistency with the PRISMA report of full-text exclusions.

\subsection{Quality Assessment}
The quality assessment was based on Section 6.3 of Kitchenham (2007) \cite{kitchenham2007}, following its recommendation to select items according to the research questions. Each checklist question was contrasted with the main review question and with directives D1--D3, retaining only those criteria that provided relevant evidence on methodological validity and robustness of results in the context of fuzzy-logic-based air-quality monitoring.

The procedure included a prior study typification. Each article included after full-text review was classified as \texttt{Qualitative}, \texttt{Quantitative}, or \texttt{Mixed}. For studies with a quantitative component, a further subclassification was made into general empirical quantitative, correlational/observational, \textit{survey}, or experiment. The checklist was applied conditionally to this typification, evaluating each study only with the items pertinent to its type and subtype.

The final instrument was built from Kitchenham's quantitative block \cite{kitchenham2007} (common questions and subtype-specific questions), applying three selection criteria: relevance to the research question, capacity to identify risk of bias (design, execution, analysis, and conclusions), and applicability to the final corpus. The result was a 30-item instrument organized into a common core and specific components for experimental, correlational/observational, and empirical quantitative studies without a specific subtype. Items associated with \textit{survey} were excluded because no studies of that type were present in the sample.

The unit of assessment was the primary study included after full-text review. Since the final corpus consisted exclusively of quantitative studies, the qualitative checklist was not applied. Coding used the scale \texttt{Yes}, \texttt{Partial}, \texttt{No}, \texttt{Unclear}, and \texttt{N/A}. The \texttt{Unclear} category was used when the reported information was insufficient to decide, and \texttt{N/A} only when the criterion did not apply to the study subtype. Quality-assessment decisions were recorded together with the supporting excerpt and the corresponding page reference (format \texttt{p.X} or \texttt{p.X-p.Y}) to preserve traceability.

The aggregated score was calculated through explicit weighting (\texttt{Yes}=1, \texttt{Partial}=0.5, \texttt{Unclear}=0.25, \texttt{No}=0), excluding \texttt{N/A} values from the denominator:

\begin{equation}
\textit{Score}_{raw} = 1\cdot n_{Yes} + 0.5\cdot n_{Partial} + 0.25\cdot n_{Unclear}
\end{equation}

\begin{equation}
\textit{Items}_{aplicables} = 30 - n_{N/A}
\end{equation}

\begin{equation}
\textit{Score}_{\%} = \left( \frac{\textit{Score}_{raw}}{\textit{Items}_{aplicables}} \right)\times 100
\end{equation}

Based on this percentage, three quality levels were defined: \textbf{High} (\(\textit{Score}_{\%}\geq 80\)), \textbf{Medium} (\(60 \leq \textit{Score}_{\%} < 80\)), and \textbf{Low} (\(\textit{Score}_{\%}<60\)). The assessment informed the interpretation of heterogeneity across findings, the identification of possible biases, and the sensitivity analysis, without being applied as an automatic exclusion criterion.

\subsection{Data Extraction}
Data extraction was defined to address the main research question through directives D1--D3 and the evaluation block, while keeping every synthesized statement linked to the corresponding full-text evidence. Following Kitchenham \cite{kitchenham2007}, the extraction instrument was specified in advance so that the required data and their coding rules were established before full extraction began.

Before full extraction, the matrix was piloted on a subset of included studies to adjust operational definitions, reduce ambiguity across fields, and verify that each variable contributed to the planned synthesis. This step resulted in a six-block matrix covering bibliographic metadata, data and variables (D1), fuzzy-system design (D2), real-time platform implementation (D3), evaluation results, and methodological traceability. The final base matrix contained 50 columns.

The matrix recorded categorical values, technical descriptions, and reported numerical results. Categorical fields were coded with a controlled catalog and explicit normalization rules. When a reported category did not fit the catalog, it was preserved as \texttt{Other:<text>}. The value \texttt{NR} was used for information not reported in the article, whereas \texttt{NA} was used when a field did not apply methodologically. This coding scheme made it possible to distinguish missing reporting from non-applicability and terminology differences.

For D1, D2, D3, and evaluation, each extracted analytical value was accompanied by page-based evidence in the format \texttt{p.X} or \texttt{p.X-p.Y}. Evidence was stored separately by block (\texttt{Evidence\_D1\_pages}, \texttt{Evidence\_D2\_pages}, \texttt{Evidence\_D3\_pages}, and \texttt{Evidence\_Eval\_pages}) so that extraction could be reviewed and audited by analytical dimension during synthesis.

The operational implementation of this phase is preserved in the final extraction matrix and in its controlled coding catalog. The conceptual view of the matrix is summarized in Table~\ref{tab:matriz_extraccion_actualizada}.

To strengthen traceability and reproducibility, the final SRL datasets and the synthesis-derived analytical artifacts used in this review were deposited in Zenodo and are publicly available at \url{https://zenodo.org/records/18929607} (DOI: \href{https://doi.org/10.5281/zenodo.18929606}{10.5281/zenodo.18929606}).

\begin{adjustwidth}{-\extralength}{0cm}
% Requiere:
% \usepackage{tabularray}
% \usepackage{caption}

\begin{longtblr}
[
 label={tab:matriz_extraccion_actualizada},
 caption={Conceptual data-extraction matrix aligned with the main question and directives D1--D3.}]{
  colspec = {X[1.2,j,m] X[3,j,m] X[3,j,m] X[0.5,j,m]},
  width = \fulllength,
  rowhead = 1,
  row{1} = {font=\bfseries}
}
\toprule
Block & Representative fields & Analytical purpose & Direc- tive \\
\midrule

Study metadata & \texttt{Article\_ID}, \texttt{Title}, \texttt{Year}, \texttt{Venue}, \texttt{DOI\_or\_URL}, \texttt{Study\_type\_quant}, \texttt{Country\_or\_region} & Identify bibliographic context, study type, and traceability of the primary source for temporal and comparative analysis. & Cross- cutting \\

Input data and variables & \texttt{D1\_Pollutants\_inputs}, \texttt{D1\_Env\_variables\_inputs}, \texttt{D1\_Target\_output}, \texttt{D1\_Sampling\_frequency\_values}, \texttt{D1\_Time\_window\_values}, \texttt{D1\_Data\_source\_type}, \texttt{D1\_Preprocessing} & Determine which pollutants and variables are used most frequently, how they are captured, and how they are prepared before modeling. & D1 \\

Fuzzy-system design & \texttt{D2\_Fuzzy\_approach\_type}, \texttt{D2\_Fuzzy\_role}, \texttt{D2\_Model\_size\_values}, \texttt{D2\_Inference\_purpose}, \texttt{D2\_Rule\_base\_definition}, \texttt{D2\_Membership\_functions}, \texttt{D2\_Defuzzification\_method}, \texttt{D2\_Optimization\_values} & Characterize the type of fuzzy inference, its centrality in the study, and the design decisions that affect interpretability and performance. & D2 \\

Platform and real-time operation & \texttt{D3\_Platform\_architecture}, \texttt{D3\_Processing\_mode}, \texttt{D3\_Communication\_and\_tech}, \texttt{D3\_Visualization\_alerting}, \texttt{D3\_Deployment\_context}, \texttt{D3\_Hardware\_and\_sensors}, \texttt{D3\_Real\_time\_characteristics} & Describe how the solution is implemented in monitoring infrastructure, its processing mode, and its actual level of continuous operation. & D3 \\

Results and performance & \texttt{Eval\_Metrics\_reported}, \texttt{Eval\_Reported\_values}, \texttt{Eval\_Comparator\_values}, \texttt{Eval\_Improvement\_reported}, \texttt{Eval\_Dataset\_size\_values}, \texttt{Eval\_Data\_split\_values}, \texttt{Eval\_Key\_results\_summary} & Compare performance across approaches, identify dominant metrics, and support analysis of effectiveness and transferability across contexts. & D1--D3 \\

Traceability and robustness for synthesis & \texttt{Main\_contribution\_reported}, \texttt{Main\_limitation\_reported}, \texttt{Evidence\_D1\_pages}, \texttt{Evidence\_D2\_pages}, \texttt{Evidence\_D3\_pages}, \texttt{Evidence\_Eval\_pages}, \texttt{QA\_Quality\_band}, \texttt{QA\_Score\_percent}, \texttt{Sensitivity\_set} & Link findings to page-verifiable evidence and integrate methodological quality into the synthesis sensitivity analysis. & Cross- cutting \\
\bottomrule
\end{longtblr}

\end{adjustwidth}


\subsubsection{Data Analysis and Synthesis}
Data analysis and synthesis were organized into five analytical phases, shown in Figure~\ref{fig:anlisis_sintesis}, in line with Kitchenham's guidance. In the first phase (analytical preparation), the consolidated extraction matrix was reviewed to verify consistency among study identifiers, analytical fields, and supporting evidence. Categorical and multi-label columns were then standardized by normalizing equivalent terms, separating compound values, and aligning terminological variants while keeping their link to the source study.

\begin{figure}[htb]
\centering
\includegraphics[width=\textwidth]{assets/figuras/metodologia/analisis_sintesis.png}
\caption{Phases of data analysis and synthesis. Own elaboration.}
\label{fig:anlisis_sintesis}
\end{figure}

In the second phase (quantitative and visual evidence), the descriptive basis was built through absolute and relative frequencies (\(k/n\), \%) and distribution tables by analytical dimension, complemented by frequency bar charts by category, distributions of the number of variables per study, and composite figures with comparative subfigures.

In the third phase (pattern identification), the analysis moved from describing categories to relating them to one another. For this purpose, cross-tabulations and co-occurrence matrices were used in order to answer how methodological decisions are combined within studies, and not only how frequent they are in isolation. At this point, convergences, recurrent configurations, and variations across approaches were identified. Descriptive columns were not forced into closed categories; instead, they were used as contextual support to interpret why certain patterns appear and how they should be read in the narrative synthesis.

In the fourth phase (profile construction), relational patterns were consolidated into compact configuration profiles. These profiles summarize typical design combinations (for example, inputs, methodological approach, or platform components) and make it possible to turn state-of-the-art findings into reusable criteria.

In the fifth phase (integration and synthesis writing), the findings were brought together into a single narrative linked to the extracted evidence and the analytical outputs. The synthesis was written to explain how the observed patterns relate to methodological decisions and design implications, rather than listing categories without interpretation.

The analytical implementation was developed in Python, mainly with \texttt{pandas} for transformation and aggregation tasks, and \texttt{seaborn}/\texttt{matplotlib} for figures and descriptive visualizations. The same workflow was applied across the synthesis directives and the evaluation blocks so that data processing, evidence tracking, and reported results followed the same analytical structure.
